% =====================================================================
\section{Conclusion}
\label{sec:conclusion}
% =====================================================================

The Euronext--TSO alliance shapes Nord Pool in five ways.

\begin{itemize}
  \item \textbf{Ecosystem coordination} between the exchange and the grid operators is governed more by regulation than by ownership. Two TSOs left the ownership structure in 2022, yet the market still functions because CACM and the shared coupling bind them to it.
  \item \textbf{Network effects:} the strongest one, market coupling, is a regulated commons shared with competitors. Nord Pool's distinctive advantage now lies in the system price, its data and the link to futures.
  \item \textbf{Platform governance} is asymmetric: Euronext owns 66\% and chairs the board, while the TSOs hold one of eight seats.
  \item \textbf{Cybersecurity} is a weakest-link problem at the interfaces, and losses fall largely on actors who make no investment decision. The pricing that grows the network also adds weak entry points.
  \item \textbf{Value creation and value capture} have diverged. Trading fees amount to roughly 0.15\% of the power price, while value capture is moving towards data and derivatives.
\end{itemize}

The most important tension revealed by our analysis is between the price as a public signal and the price as a private asset. A listed exchange group controls the platform that produces the Nordic reference price. The state-owned TSOs supply its critical input and carry the system risk, but have limited influence over how the platform is run. The 2022 data paywall \citep{ERR2022} and the 2026 launch of futures based on the system price \citep{Euronext2026} show this tension in practice. Value created jointly by members, TSOs and regulators is increasingly captured through products controlled by one partner.

This tension does not mean that the alliance destroys value. Euronext has brought capital and technology, and it took over the Nordic power futures market from Nasdaq. Multi-NEMO competition also keeps trading fees low. But it does shift the burden of protecting the public interest from ownership to regulation. Whether the alliance remains legitimate will depend on three developments:
\begin{enumerate}
  \item whether regulators treat price data as public infrastructure;
  \item whether competition from EPEX remains credible;
  \item whether the TSOs, who are indispensable but increasingly peripheral as owners, keep enough influence to protect security of supply.
\end{enumerate}
These questions link directly to the broader economic and societal assessment in our essay.
